% =============================================================
%  第六单元  图形的运动（一）
% =============================================================
\Chapter{第六关\quad 剪一刀，动一下}
\noindent{\small\kaishu\color{tipc}（课本第六单元：图形的运动（一），P72--P76）}

\begin{blunder}
过年想露一手剪窗花。剪完展开，小豆角看了半天说：
“豆角爸爸，这剪的是怪兽吗？”——确实像个怪兽，一边大一边小。
原因很简单：豆角爸爸折的时候没对齐，还从开口那边下剪子。
\end{blunder}

\begin{gamble}{赌一把：谁剪得更像}
\begin{itemize}
  \item 同一个图案（蝴蝶、树、瓶子），两人各剪一个。
  \item \textbf{请豆角妈妈来当评委}（或者拍照发给奶奶），她说了算。
  \item 输的人负责把纸屑扫干净——剪窗花最大的成本是扫地。
\end{itemize}
\end{gamble}

\begin{bonus}
\begin{itemize}
  \item \textbf{剪一串手拉手小人}贴窗户。诀窍：纸折三层，画半个人，
        手要画在折痕那一侧。剪成了就贴满阳台窗户，小豆角会有成就感。
  \item \textbf{洗澡时玩对称}：对着镜子做动作，问“镜子里的小豆角举的是左手还是右手？”
        （其实镜像反过来了，小豆角会想很久。）
  \item \textbf{找对称的招牌}：出门路上找五个左右对称的招牌或建筑，
        用手比划对称轴在哪。
\end{itemize}
\end{bonus}

\begin{unitinfo}
\textbf{课本 P72--P76}\quad \textbf{3 次课，每次 20--25 分钟}\quad
\textbf{建议第 17 周} \\[4pt]
\textbf{这一关到底在干嘛}：剪对称的纸、移动积木、转陀螺。
听起来像手工课，但考点很实在：\textbf{判断一个图形对不对称}、
\textbf{说清往哪个方向移了几格}。\\
\textbf{逻辑是这么走的}：\ 折一下 $\to$ 两边一样 $\to$ 这叫对称
$\to$ 只挪位置不转向，叫平移 $\to$ 绕着一个点转圈，叫旋转。\\
\textbf{这一章剪二十分钟纸，顶做二十道题。}别跳。
\end{unitinfo}

\subsection{开局先聊个事：这三件事每天都在发生}

小豆角穿衣服，扣子左右对称；剪纸剪窗花，一张纸对折一下剪；
推拉门只能左右拉——这叫平移；风扇、车轮、陀螺——这叫旋转。

课本里还有两个老玩具：\textbf{华容道}（三国故事）和\textbf{竹蜻蜓}。
这些不是摆设，都是让小豆角用手感受“东西怎么动”。

\subsection{任务 1：折一下，就对称了}

\begin{lesson}{1}{25}
\textbf{今天要拿下的}：会“先想怎么折、再想怎么画、最后才剪”，
剪出一个对称图形，并说清它为什么对称。

\begin{materials}
正方形手工纸 15\,cm（至少 10 张）、剪刀、铅笔、旧报纸（练手）。
\end{materials}

\begin{enumerate}
  \item \textbf{先给结果，让小豆角倒推}（课本 P72 的思路）：
        豆角爸爸拿出一个提前剪好的“手拉手小人”，问：“小豆角猜豆角爸爸怎么剪的？”
        小豆角多半说“画一个剪一个”。\textbf{让小豆角试}——剪出来只有一个人。
  \item \textbf{给个提示}：“豆角爸爸剪之前只画了\textbf{一半}。”
        让小豆角自己想到“纸要折起来”。
  \item \textbf{正式开始，三步，每步先说出来再做}：
        \begin{enumerate}[label=第\arabic*步：]
          \item \textbf{折}：对折，\textbf{边和边必须完全对齐}，用指甲把折痕刮清。
                \textbf{对不齐，剪出来就是歪的}——第一个坎。
          \item \textbf{画}：在折好的一面画半个图案。\textbf{关键：}
                图案的“身体”要画在\textbf{折痕那一侧}，开口朝外。
                画错了，展开就断了。
          \item \textbf{剪}：沿线剪，\textbf{折痕边上不能剪断}（那里是连接处）。
        \end{enumerate}
  \item \textbf{展开}：让小豆角自己展开，看到“两个人手拉手”。
        问：“为什么变成两个了？”（折了两层，剪一次等于剪两次）
  \item \textbf{再剪三个不同的}：树、花瓶、蝴蝶。
        每剪一个都问“小豆角把什么画在折痕上？”
  \item \textbf{判断对称（考点）}：给几个图形，判断哪些对称，
        并\textbf{用手指比划折痕在哪}。
        方法就一句：\textbf{能不能找到一条线，沿它折过去两边完全重合}。
  \item \textbf{找家里的对称}：窗、门、杯子、蝴蝶、汉字“中”“田”“喜”。
        再摸自己的脸——“两边也是一样的”。
\end{enumerate}

\begin{saybox}
“小豆角只画了一半，为什么能剪出两个？”\\
“折的时候没对齐，剪出来会怎么样？”（让小豆角自己试一次错的）\\
“折痕要是剪断了，会变成什么？”\\
“小豆角能在屋里找出三个对称的东西吗？”
\end{saybox}

\begin{pitbox}
\textbf{三个坑，建议让小豆角先踩一次再纠}：
\begin{enumerate}[label=\arabic*、]
  \item 折痕没对齐 $\to$ 剪出来上下不一样大。
  \item 图案画在开口那一边 $\to$ 展开是两个断开的碎片。
  \item 沿折痕剪断了 $\to$ 变成两块单独的图形，不连着了。
\end{enumerate}
\textbf{别提前全讲一遍}，让小豆角先错，错了问小豆角：“小豆角觉得是哪一步出的问题？”
\end{pitbox}

\begin{passbox}
\begin{itemize}
  \item 独立剪出一个对称图形，能说出“折—画—剪”三步。
  \item 能判断 4 个图形是否对称，指出折痕位置。
  \item 能说出“剪之前要想清楚怎么折”。
\end{itemize}
\end{passbox}
\end{lesson}

\subsection{任务 2：华容道——只能横着走或竖着走}

\begin{lesson}{2}{20}
\textbf{今天要拿下的}：理解平移就是只换位置、不换方向；能说清“往哪边移了几格”。

\begin{materials}
快递纸箱（做底板）、卡纸剪的人物、铅笔、尺子。
\end{materials}

\begin{enumerate}
  \item \textbf{做棋盘（提前五分钟）}：纸箱上画 4 列 $\times$ 5 行的格子，
        右下角留两格出口。卡纸剪四块人物牌——
        曹操（2$\times$2）、关羽（2$\times$1）、赵云（1$\times$2）、小兵（1$\times$1 两块）。
  \item \textbf{讲规则，只说三条}：
        ① 所有人物\textbf{只能横着或竖着移}，不能斜着、不能提起来；
        ② 不能穿过别人；
        ③ 曹操移到出口就赢。
  \item \textbf{先让小豆角乱玩三分钟}：小豆角会发现“走不动”“卡住了”。
        \textbf{这个“卡住”就是今天要教的东西。}
  \item \textbf{教“说清楚”}：不许比划，必须说一句话——
        “\textbf{把赵云向上移动 3 格。}”\\
        让小豆角每次移动前\textbf{先说出方向和格数}，豆角爸爸照着执行（小豆角指挥豆角爸爸）。
        \textbf{这一步才是考点训练。}
  \item \textbf{数格子}：横着走 2 格、竖着走 3 格，要说清“从哪一格到哪一格”。
  \item \textbf{记步数}：用“正”字记，两轮比比谁步数少。
  \item \textbf{换过来}：豆角爸爸走一步，让小豆角说“小豆角把谁往哪边移了几格”。
\end{enumerate}

\begin{saybox}
“小豆角要把它移到哪儿？上还是左？几格？”（每步都逼小豆角说完整）\\
“为什么这一步走不动？”（因为被挡住了——这就是理解规则）\\
“刚才那一步，能不能换个走法？”
\end{saybox}

\begin{passbox}
\begin{itemize}
  \item 能用“向（上／下／左／右）移动（几）格”准确说清 5 次移动。
  \item 能说出为什么某步走不通。
\end{itemize}
\end{passbox}
\end{lesson}

\subsection{任务 3：陀螺、风车、竹蜻蜓}

\begin{lesson}{3}{25}
\textbf{今天要拿下的}：在玩里体会“旋转”；自己发现“轴扎在中间最稳”。

\begin{materials}
卡纸、竹签、图钉、吸管、剪刀、尺子。
\end{materials}

\begin{enumerate}
  \item \textbf{做三个陀螺}（课本 P75）：剪圆形、正方形、三角形各一个，掌心大小。
        \begin{itemize}
          \item 圆形：竹签扎在\textbf{正中心}（用尺子找：两条直径的交点）。
          \item 正方形：扎在\textbf{对角线交点}。
          \item 三角形：\textbf{故意扎偏一点}，做成“偏心”的。
        \end{itemize}
  \item \textbf{比一比}（今天的核心发现）：三个各转三次，记录稳不稳。
        \begin{center}
        \small
        \begin{tabular}{@{}p{3.8cm}p{2.8cm}p{2.8cm}p{3.0cm}@{}}
        \toprule
        \textbf{陀螺} & \textbf{第一次} & \textbf{第二次} & \textbf{结论} \\
        \midrule
        圆形（扎中心） & 稳／不稳 & 稳／不稳 & \blank[2.0cm] \\
        正方形（扎中心） & 稳／不稳 & 稳／不稳 & \blank[2.0cm] \\
        三角形（扎偏心） & 稳／不稳 & 稳／不稳 & \blank[2.0cm] \\
        \bottomrule
        \end{tabular}
        \end{center}
  \item \textbf{让小豆角自己下结论}：“\textbf{扎在中间最稳，扎偏了就晃。}”
        \textbf{豆角爸爸先不说}，让小豆角从记录里看出来。
  \item \textbf{再验证一次（最重要的一步）}：把圆形那个的竹签故意挪偏 1 厘米，再转。
        小豆角会看到“同一个陀螺，扎偏了也不稳”。
        \textbf{这一步控制了变量}，比讲一百遍“中心”有用。
  \item \textbf{做风车}（课本 P76）：正方形纸四角向中心剪（\textbf{别剪到中心，留 1 厘米}），
        四个角交错折向中心，图钉从中心扎过去，固定在铅笔或吸管上。
  \item \textbf{做竹蜻蜓}：卡纸剪长条当翅膀，插进吸管，用胶或胶带固定。两手一搓，飞起来。
  \item \textbf{收口}：让小豆角说这三个玩具的共同点——
        \textbf{都是绕着一个点转圈}，这就是\textbf{旋转}。
\end{enumerate}

\begin{saybox}
“哪个转得最稳？为什么？”\\
“同一个圆陀螺，扎中间和扎偏，差在哪？”\\
“陀螺、风车、竹蜻蜓，转的时候都绕着哪儿转？”
\end{saybox}

\begin{tipbox}
\textbf{安全}：竹签和图钉都要豆角爸爸在旁边看着。陀螺在桌面上转，别举高。
剪刀用完立刻收走。
\end{tipbox}

\begin{passbox}
\begin{itemize}
  \item 三个陀螺都做出来、都转起来了。
  \item 小豆角能从记录表里自己说出“扎中心最稳”。
  \item 小豆角能说出三样玩具的共同点（绕一个点转）。
\end{itemize}
\end{passbox}
\end{lesson}

\subsection{这一章的三个词，说清就够}

\begin{center}
\small
\begin{tabularx}{\linewidth}{@{}p{2.2cm}p{4.4cm}X@{}}
\toprule
\textbf{词} & \textbf{课本里怎么玩} & \textbf{小豆角要会说的一句话} \\
\midrule
对称 & 对折剪小人、树、花瓶 & “对折后两边一样，折痕就是对称轴。” \\
平移 & 华容道 & “只能横着或竖着移，方向不变，移了 3 格。” \\
旋转 & 陀螺、风车、竹蜻蜓 & “绕着一个点转圈。” \\
\bottomrule
\end{tabularx}
\end{center}

\noindent\textbf{生活里接着练}：逛街找对称的建筑和招牌；坐电梯看“上升”；
玩滑板车问“轮子是平移还是旋转？”（两个都有）。
